\chapter{การจำลองฟิสิกส์เชิงสะเต็มในโลกเสมือนจริง}
\label{chap:ch04}

\section*{แผนบริหารการสอนประจำบทที่ 4}
\addcontentsline{toc}{section}{แผนบริหารการสอนประจำบทที่ 4}

\noindent\textbf{หัวข้อเนื้อหาประจำบท}
\begin{enumerate}
    \item ระเบียบวิธีเชิงตัวเลขสำหรับการจำลองการเคลื่อนที่ (Euler, Verlet, and RK4)
    \item พลศาสตร์ของวัตถุแข็งเกร็ง แรงบิด และโมเมนต์ความเฉื่อยใน 3 มิติ
    \item ขั้นตอนวิธีตรวจจับการชนและการคำนวณการดลเพื่อตอบสนองการกระดอน
    \item การจำลองระบบการสั่น ระบบสปริงหน่วง และสนามความโน้มถ่วงในโลกเสมือน
\end{enumerate}

\noindent\textbf{วัตถุประสงค์เชิงพฤติกรรม}
\begin{enumerate}
    \item เปรียบเทียบความแม่นยำและความเสถียรของวิธี Euler, Verlet, และ RK4 ได้
    \item ประยุกต์กฎการเคลื่อนที่ของนิวตันในการอัปเดตตำแหน่งและความเร็วของวัตถุได้
    \item คำนวณการชนแบบยืดหยุ่นและไม่ยืดหยุ่นผ่านสมการการดล (Impulse Response) ได้
    \item พัฒนาเอนจินจำลองลูกตุ้มเพนดูลัมและสปริงที่สอดคล้องกับทฤษฎีทางฟิสิกส์ได้
\end{enumerate}

\newpage

\noindent\textbf{กิจกรรมการเรียนการสอน}
\begin{enumerate}
    \item การบรรยายเชิงคณิตศาสตร์ฟิสิกส์และการสาธิตการสะสมความคลาดเคลื่อนเชิงตัวเลข
    \item กิจกรรมการเขียนโค้ดเปรียบเทียบวิถีโค้งของวัตถุภายใต้สนามโน้มถ่วง
    \item การทดลองจำลองการชนของมวลในระบบหนึ่งมิติและสองมิติในสภาพแวดล้อมเสมือน
    \item การวิเคราะห์ข้อมูลการสั่นแบบหน่วงเปรียบเทียบระหว่างแบบจำลองกับผลการคำนวณวิเคราะห์
\end{enumerate}

\noindent\textbf{สื่อการเรียนการสอน}
\begin{enumerate}
    \item เอกสารประกอบการสอนและสไลด์สรุปสูตรฟิสิกส์เชิงคำนวณ
    \item โค้ดจำลอง Physics Loop ในภาษา Python (NumPy) และ WebAssembly
    \item กราฟแสดงการอนุรักษ์พลังงานกลรวมในระบบสปริงแบบเรียลไทม์
    \item แอปพลิเคชันจำลองการทดลองเสมือนจริงบนมาตรฐาน OpenXR
\end{enumerate}

\noindent\textbf{การประเมินผล}
\begin{enumerate}
    \item การประเมินความถูกต้องของแบบจำลองฟิสิกส์ที่ผู้เรียนพัฒนา
    \item การตรวจรายงานการวิเคราะห์พลังงานและเสถียรภาพของตัวอินทิเกรตเชิงตัวเลข
    \item การทดสอบข้อเขียนเชิงทฤษฎีและปฏิบัติการคำนวณ
\end{enumerate}

\newpage

\section{ระเบียบวิธีเชิงตัวเลขสำหรับการจำลองการเคลื่อนที่}

หัวใจสำคัญของห้องปฏิบัติการเสมือนจริงทางสะเต็มศึกษาคือความถูกต้องแม่นยำของกฎฟิสิกส์ที่ควบคุมวัตถุ ในคอมพิวเตอร์ การแก้สมการเชิงอนุพันธ์ของการเคลื่อนที่ $\mathbf{F} = m\frac{d^2\mathbf{x}}{dt^2}$ ต้องอาศัยระเบียบวิธีเชิงตัวเลข (Numerical Integrators)

\subsection{วิธีออยเลอร์แบบดั้งเดิมและแบบกึ่งชัดแจ้ง}

วิธีออยเลอร์แบบชัดแจ้ง (Explicit Euler) ทำการประมาณค่าตำแหน่งและความเร็วในก้าวเวลา $\Delta t$ ถัดไป ดังนี้
\begin{align}
\mathbf{x}_{n+1} &= \mathbf{x}_n + \mathbf{v}_n \Delta t \\
\mathbf{v}_{n+1} &= \mathbf{v}_n + \mathbf{a}_n \Delta t
\end{align}
แม้วิธีนี้จะคำนวณง่าย แต่มีข้อเสียร้ายแรงคือพลังงานของระบบจะไม่ได้รับการอนุรักษ์ และมักเกิดการระเบิดเชิงตัวเลข (Numerical Explosion) ในระบบที่มีการสั่น จึงนิยมใช้วิธีออยเลอร์กึ่งชัดแจ้ง (Semi-Implicit Euler หรือ Symplectic Euler) แทน
\begin{align}
\mathbf{v}_{n+1} &= \mathbf{v}_n + \mathbf{a}_n \Delta t \\
\mathbf{x}_{n+1} &= \mathbf{x}_n + \mathbf{v}_{n+1} \Delta t
\end{align}

\subsection{วิธีแวร์เลต์และความเร็วแวร์เลต์}

วิธีแวร์เลต์ (Verlet Integration) เป็นระเบียบวิธีที่นิยมใช้อย่างแพร่หลายในเอนจินฟิสิกส์ เนื่องจากมีความเสถียรสูงและอนุรักษ์พลังงานในระบบอนุรักษ์ได้อย่างยอดเยี่ยม
\begin{equation}
\mathbf{x}_{n+1} = 2\mathbf{x}_n - \mathbf{x}_{n-1} + \mathbf{a}_n \Delta t^2
\end{equation}
ส่วนวิธี Velocity Verlet ช่วยให้สามารถคำนวณความเร็วพร้อมกันได้ในแต่ละก้าวเวลา
\begin{align}
\mathbf{x}_{n+1} &= \mathbf{x}_n + \mathbf{v}_n \Delta t + \frac{1}{2}\mathbf{a}_n \Delta t^2 \\
\mathbf{v}_{n+1} &= \mathbf{v}_n + \frac{1}{2}(\mathbf{a}_n + \mathbf{a}_{n+1})\Delta t
\end{align}

\begin{figure}[htbp]
\centering
\begin{tikzpicture}[scale=1.1]
    % Axes
    \draw[-{Stealth}, line width=1.2pt] (0,0) -- (6.5,0) node[right]{เวลา $t$};
    \draw[-{Stealth}, line width=1.2pt] (0,-2.2) -- (0,2.5) node[above]{การกระจัด $x(t)$};

    % Analytical exact curve
    \draw[darkNavy, line width=1.4pt, domain=0:6.0, samples=100] plot (\x, {1.8*cos(2.0*\x r)});
    \node[darkNavy, font=\small] at (5.0, 2.2) {\textbf{ผลเฉลยวิเคราะห์ (Exact)}};

    % Explicit Euler Divergence
    \draw[red!80!black, dashed, line width=1.3pt, domain=0:5.5, samples=12] plot (\x, {1.8*exp(0.12*\x)*cos(2.0*\x r)});
    \node[red!80!black, font=\small] at (5.5, -1.8) {\textbf{Explicit Euler (แกว่งเกินและสูญเสียเสถียรภาพ)}};

    % Symplectic / Verlet
    \draw[emeraldActive, dotted, line width=1.5pt, domain=0:6.0, samples=20] plot (\x, {1.8*cos(2.02*\x r)});
    \node[emeraldActive, font=\small] at (2.5, -2.4) {\textbf{Verlet / Symplectic (อนุรักษ์พลังงาน)}};
\end{tikzpicture}
\caption{การเปรียบเทียบเสถียรภาพเชิงตัวเลขระหว่างวิธี Explicit Euler และวิธี Verlet ในระบบการสั่นฮาร์มอนิกอย่างง่าย}
\label{fig:numerical_stability}
\end{figure}

\section{การตรวจจับการชนและการคำนวณแรงดล}

ในการจำลองสะเต็ม วัตถุเสมือนต้องสามารถชนและสะท้อนกลับตามกฎการอนุรักษ์โมเมนตัม การชนระหว่างวัตถุสองชิ้นที่มีมวล $m_1, m_2$ และความเร็วก่อนชน $\mathbf{v}_1, \mathbf{v}_2$ สามารถคำนวณได้จากขนาดของการดล $J$ ตามแนวตั้งฉากของการชน $\mathbf{n}$

\begin{equation}
J = \frac{-(1 + e)(\mathbf{v}_1 - \mathbf{v}_2)\cdot \mathbf{n}}{\frac{1}{m_1} + \frac{1}{m_2}}
\end{equation}
โดยที่ $e \in [0, 1]$ คือสัมประสิทธิ์การคืนตัว (Coefficient of Restitution) ความเร็วหลังการชนคำนวณได้จาก
\begin{align}
\mathbf{v}_1' &= \mathbf{v}_1 + \frac{J}{m_1}\mathbf{n} \\
\mathbf{v}_2' &= \mathbf{v}_2 - \frac{J}{m_2}\mathbf{n}
\end{align}

\begin{workedexamplebox}{การคำนวณการชนยืดหยุ่นของลูกกลมโลหะสองลูกบนรางเอียงเสมือน}
ลูกกลมโลหะ $A$ มวล $m_A = 0.2\text{ kg}$ เคลื่อนที่ด้วยความเร็ว $\mathbf{v}_A = (1.5, 0, 0)\text{ m/s}$ เข้าชนลูกกลม $B$ มวล $m_B = 0.3\text{ kg}$ ซึ่งหยุดนิ่งอยู่กับที่ ($\mathbf{v}_B = (0, 0, 0)\text{ m/s}$) หากการชนเป็นการชนแบบยืดหยุ่นสมบูรณ์ ($e = 1.0$) โดยเวกเตอร์แนวตั้งฉากของการชนคือ $\mathbf{n} = (1, 0, 0)$ จงหาความเร็วของลูกกลมทั้งสองหลังการชน

\textbf{วิธีทำ:}
คำนวณความเร็วสัมพัทธ์ในแนวตั้งฉาก
\begin{equation}
(\mathbf{v}_A - \mathbf{v}_B)\cdot \mathbf{n} = (1.5 - 0)(1) = 1.5\text{ m/s}
\end{equation}
คำนวณผลรวมส่วนกลับของมวล
\begin{equation}
\frac{1}{m_A} + \frac{1}{m_B} = \frac{1}{0.2} + \frac{1}{0.3} = 5.0 + 3.333 = 8.333\text{ kg}^{-1}
\end{equation}
คำนวณขนาดของการดล $J$
\begin{equation}
J = \frac{-(1 + 1.0)(1.5)}{8.333} = \frac{-3.0}{8.333} \approx -0.360\text{ N}\cdot\text{s}
\end{equation}
คำนวณความเร็วหลังการชนของลูกกลมแต่ละลูก
\begin{align}
\mathbf{v}_A' &= (1.5, 0, 0) + \frac{-0.360}{0.2}(1, 0, 0) = (1.5 - 1.80, 0, 0) = (-0.30, 0, 0)\text{ m/s} \\
\mathbf{v}_B' &= (0, 0, 0) - \frac{-0.360}{0.3}(1, 0, 0) = (0 + 1.20, 0, 0) = (1.20, 0, 0)\text{ m/s}
\end{align}
ผลลัพธ์แสดงว่าลูกกลม $A$ สะท้อนถอยหลังด้วยความเร็ว $0.30\text{ m/s}$ ขณะที่ลูกกลม $B$ เคลื่อนที่ไปข้างหน้าด้วยความเร็ว $1.20\text{ m/s}$ ตรงตามกฎการอนุรักษ์โมเมนตัมและพลังงานกล
\end{workedexamplebox}

\section{การจำลองระบบสปริงหน่วงและการสั่นพ้อง}

ระบบมวล-สปริงแบบมีตัวหน่วง (Damped Harmonic Oscillator) สอดคล้องกับสมการอนุพันธ์อันดับสอง
\begin{equation}
m\frac{d^2 x}{dt^2} + c\frac{dx}{dt} + kx = F_{\text{ext}}(t)
\end{equation}
โดยที่ $k$ คือค่านิจสปริง และ $c$ คือสัมประสิทธิ์ความหน่วง ในการจำลองเชิงสะเต็ม ผู้เรียนสามารถปรับเปลี่ยนอัตราส่วนความหน่วง $\zeta = \frac{c}{2\sqrt{mk}}$ เพื่อสังเกตปรากฏการณ์กวัดแกว่งแบบต่าง ๆ ได้แก่ สั่นหน่วงน้อย (Underdamped, $\zeta < 1$), สั่นหน่วงวิกฤต (Critically Damped, $\zeta = 1$), และสั่นหน่วงเกิน (Overdamped, $\zeta > 1$) ดังแสดงในตารางที่ 4.1

\begin{table}[htbp]
\centering
\caption{สภาวะการเคลื่อนที่ของระบบมวลสปริงแบบมีตัวหน่วงในห้องปฏิบัติการเสมือน}
\label{tab:damping_states}
\begin{tabularx}{\textwidth}{lYY}
\toprule
\textbf{สภาวะความหน่วง} & \textbf{เงื่อนไขพารามิเตอร์} & \textbf{พฤติกรรมทางกายภาพที่สังเกตได้} \\
\midrule
หน่วงน้อย (Underdamped) & $\zeta < 1.0$ & มีการแกว่งสลับไปมาโดยแอมพลิจูดลดลงแบบเอกซ์โพเนนเชียล \\
หน่วงวิกฤต (Critical) & $\zeta = 1.0$ & มวลกลับคืนสู่จุดสมดุลเร็วที่สุดโดยไม่มีการแกว่งข้ามศูนย์ \\
หน่วงเกิน (Overdamped) & $\zeta > 1.0$ & มวลเคลื่อนที่เข้าสู่จุดสมดุลอย่างช้า ๆ เนื่องจากแรงต้านทานสูงมาก \\
\bottomrule
\end{tabularx}
\end{table}

\begin{aipromptbox}[การสร้างเอนจินฟิสิกส์ระบบสปริงหน่วงด้วยระเบียบวิธีเชิงตัวเลขด้วย AI]
\textbf{วัตถุประสงค์:} พัฒนาโปรแกรมจำลองฟิสิกส์สำหรับระบบมวล-สปริงแบบมีตัวหน่วง (Damped Harmonic Oscillator) ที่มีความแม่นยำสูงด้วยระเบียบวิธี Velocity Verlet พร้อมวิเคราะห์การอนุรักษ์พลังงานกลรวม\\
\textbf{โครงสร้าง CLEAR Prompt:}
\begin{itemize}[topsep=1pt, itemsep=1pt, leftmargin=1.2em]
    \item \textbf{Context (บริบท):} โมดูลการทดลองฟิสิกส์เสมือนจริงบน OpenXR ต้องการคำนวณการเคลื่อนที่ของมวลและแรงดึงกลับแบบเรียลไทม์
    \item \textbf{Logic (ตรรกะ):} แก้สมการ $m\ddot{x} + c\dot{x} + kx = 0$ ด้วย Velocity Verlet อัปเดต $x(t+\Delta t)$ และ $v(t+\Delta t)$ พร้อมคำนวณพลังงานจลน์ $E_k = \frac{1}{2}mv^2$ และพลังงานศักย์ยืดหยุ่น $E_p = \frac{1}{2}kx^2$
    \item \textbf{Expectation (ผลลัพธ์ที่ต้องการ):} คลาส Python \texttt{SpringDamperSimulator} ที่มีเมท็อด \texttt{step(dt)} และคืนค่าสถานะพิกัดพร้อมพลังงานรวม
    \item \textbf{Action \& Restriction (ข้อกำหนด):} ต้องรองรับการเปลี่ยนค่าคงที่สปริง $k$ และสัมประสิทธิ์หน่วง $c$ ขณะรัน และมีตัวอย่างโค้ดรันแสดงผล 5 ก้าวเวลา
\end{itemize}
\end{aipromptbox}

\begin{codebox}[การจำลองระบบมวลสปริงหน่วงด้วยวิธี Velocity Verlet ในภาษา Python]
import numpy as np

class SpringDamperSimulator:
    def __init__(self, mass=1.0, k=50.0, c=0.5, x0=0.2, v0=0.0):
        self.m = mass
        self.k = k
        self.c = c
        self.x = x0
        self.v = v0
        self.a = self._compute_accel(self.x, self.v)

    def _compute_accel(self, x, v):
        # F = -kx - cv => a = F/m
        return (-self.k * x - self.c * v) / self.m

    def step(self, dt):
        # ขั้นที่ 1: อัปเดตตำแหน่งด้วย Velocity Verlet
        self.x += self.v * dt + 0.5 * self.a * (dt ** 2)
        
        # ขั้นที่ 2: ประมาณความเร็วกึ่งก้าวเพื่อหาความเร่งใหม่
        v_half = self.v + 0.5 * self.a * dt
        a_new = self._compute_accel(self.x, v_half)
        
        # ขั้นที่ 3: อัปเดตความเร็วเต็มก้าว
        self.v = v_half + 0.5 * a_new * dt
        self.a = a_new
        
        # คำนวณพลังงานรวม (Mechanical Energy)
        e_kinetic = 0.5 * self.m * (self.v ** 2)
        e_potential = 0.5 * self.k * (self.x ** 2)
        return self.x, self.v, e_kinetic + e_potential

# ตัวอย่างการทำงานที่ก้าวเวลา dt = 0.01 วินาที (100 Hz)
sim = SpringDamperSimulator(mass=0.5, k=20.0, c=0.2, x0=0.1)
print(f"Initial: x={sim.x:.4f} m, v={sim.v:.4f} m/s")
for step_idx in range(1, 6):
    x, v, e_total = sim.step(dt=0.01)
    print(f"Step {step_idx}: x={x:.4f} m, v={v:.4f} m/s, E_total={e_total:.5f} J")
\end{codebox}

\section{ปฏิบัติการสะเต็มฟิสิกส์เสมือนจริง: ลูกตุ้ม เพนดูลัม และพื้นเอียง}

เพื่อเชื่อมโยงทฤษฎีเชิงตัวเลขเข้ากับการทดลองวิทยาศาสตร์จริง ห้องปฏิบัติการเสมือนจริงบนสถาปัตยกรรม OpenXR สามารถจำลองชุดการทดลองฟิสิกส์มาตรฐานที่ผู้เรียนสามารถโต้ตอบ ปรับค่าพารามิเตอร์ และบันทึกผลได้แบบเรียลไทม์

\subsection{การจำลองลูกตุ้มอย่างง่ายและการอนุรักษ์พลังงานในระนาบเฟส}

การเคลื่อนที่ของลูกตุ้มอย่างง่าย (Simple Pendulum) มวล $m$ แขวนด้วยเชือกยาว $L$ สอดคล้องกับสมการไม่เชิงเส้น
\begin{equation}
\frac{d^2\theta}{dt^2} + \frac{g}{L}\sin\theta = 0 \label{eq:simple_pendulum}
\end{equation}
เมื่อมุมเบี่ยงเบน $\theta$ มีค่าน้อย ($\sin\theta \approx \theta$) ระบบจะประมาณเป็นการสั่นแบบฮาร์มอนิกอย่างง่ายที่มีคาบการแกว่ง $T = 2\pi\sqrt{\frac{L}{g}}$ แต่ในห้องปฏิบัติการเสมือน ผู้เรียนสามารถดึงลูกตุ้มด้วยมือจนมีมุมกว้าง ($\theta > 30^\circ$) จึงต้องแก้สมการ \eqref{eq:simple_pendulum} ด้วยระเบียบวิธีเชิงตัวเลขแบบเรียลไทม์

พลังงานกลรวมของลูกตุ้มประกอบด้วยพลังงานจลน์ $E_k$ และพลังงานศักย์โน้มถ่วง $E_p$
\begin{align}
E_k &= \frac{1}{2} m v^2 = \frac{1}{2} m (L\dot{\theta})^2 \label{eq:pendulum_ek} \\
E_p &= mgL(1 - \cos\theta) \label{eq:pendulum_ep} \\
E_{\text{total}} &= E_k + E_p = \text{คงที่} \label{eq:pendulum_etotal}
\end{align}
การแสดงผลกราฟพลังงานควบคู่กับระนาบเฟส (Phase Portrait: $\theta$ เทียบกับ $\dot{\theta}$) ช่วยให้ผู้เรียนมองเห็นวงโคจรแบบปิด (Closed Orbit) อันเป็นเอกลักษณ์ของระบบอนุรักษ์

\subsection{ปฏิบัติการแรงเสียดทานและพื้นเอียงปรับมุมด้วยมือ}

การทดลองหาค่าสัมประสิทธิ์แรงเสียดทานด้วยพื้นเอียงเสมือนจริง ผู้เรียนสามารถใช้ท่าทางมือ (Hand Gesture) แตะขอบพื้นเอียงแล้วยกปรับมุมเท $\alpha$ ได้อย่างอิสระ วัตถุมวล $m$ ที่วางอยู่บนพื้นเอียงจะอยู่ภายใต้แรงกระทำ 3 แรงหลัก:
\begin{enumerate}
    \item แรงโน้มถ่วงในแนวดิ่ง $W = mg$ ซึ่งแยกเป็นองค์ประกอบตั้งฉาก $mg\cos\alpha$ และขนานพื้นเอียง $mg\sin\alpha$
    \item แรงแนวฉากที่พื้นกระทำต่อวัตถุ $N = mg\cos\alpha$
    \item แรงเสียดทาน $f$ ต้านการเลื่อนไถล
\end{enumerate}

\begin{figure}[htbp]
\centering
\begin{tikzpicture}[scale=1.1]
    % Inclined Plane
    \draw[line width=1.5pt, darkNavy] (0,0) -- (6,0) -- (6,3.0) -- cycle;
    \draw[line width=1.2pt, slateGrey] (0,0) -- (6,3.0);
    \node at (1.5,0.3) {$\alpha$};
    \draw[->, cyberPurple, line width=1.0pt] (1.8,0) arc (0:26.5:1.8);

    % Block on Plane
    \begin{scope}[shift={(3.5, 1.75)}, rotate=26.5]
        \draw[fill=openxrTeal!20, draw=openxrTeal, line width=1.4pt] (-0.7,0) rectangle (0.7,0.8);
        \fill[darkNavy] (0,0.4) circle (2.5pt) node[above right]{\textbf{$m$}};

        % Normal force N
        \draw[-{Stealth}, line width=1.5pt, emeraldActive] (0,0.4) -- (0, 2.0) node[above]{\textbf{$\mathbf{N}$}};

        % Friction force f
        \draw[-{Stealth}, line width=1.5pt, amberGlow] (0,0.4) -- (-1.6, 0.4) node[left]{\textbf{$\mathbf{f}$}};

        % Parallel gravity
        \draw[-{Stealth}, line width=1.2pt, red!80!black] (0,0.4) -- (1.6, 0.4) node[right]{\textbf{$mg\sin\alpha$}};
    \end{scope}

    % Gravity vector downwards
    \draw[-{Stealth}, line width=1.5pt, spatialBlue] (3.5, 2.1) -- (3.5, 0.2) node[below]{\textbf{$mg$}};
\end{tikzpicture}
\caption{แผนภาพแรงเสรี (Free-Body Diagram) ของวัตถุบนพื้นเอียงปรับมุมในห้องปฏิบัติการเสมือน}
\label{fig:inclined_plane}
\end{figure}

ขณะที่วัตถุยังอยู่นิ่ง แรงเสียดทานสถิตจะสมดุลกับแรงโน้มถ่วง $f_s = mg\sin\alpha$ จนกระทั่งมุมเอียงเพิ่มขึ้นถึง \textbf{มุมวิกฤต (Angle of Repose)} $\alpha_c$ ซึ่งเป็นมุมที่วัตถุเริ่มไถล แรงเสียดทานสถิตจะมีค่าสูงสุด $f_{s,\max} = \mu_s N$
\begin{equation}
mg\sin\alpha_c = \mu_s mg\cos\alpha_c \implies \mu_s = \tan\alpha_c \label{eq:repose_angle}
\end{equation}
เมื่อมุมเอียง $\alpha > \alpha_c$ วัตถุจะไถลลงด้วยความเร่ง $a$ ภายใต้แรงเสียดทานจลน์ $f_k = \mu_k N$
\begin{equation}
a = g(\sin\alpha - \mu_k \cos\alpha) \label{eq:kinetic_accel}
\end{equation}

\begin{workedexamplebox}{การคำนวณสัมประสิทธิ์แรงเสียดทานสถิตและความเร่งบนพื้นเอียง}
ในการทดลองเสมือนจริง ผู้เรียนปรับมุมของพื้นเอียงจนพบว่าบล็อกไม้เริ่มไถลลงที่มุม $\alpha_c = 30^\circ$ หลังจากนั้นผู้เรียนปรับมุมเพิ่มขึ้นเป็น $\alpha = 45^\circ$ หากสัมประสิทธิ์แรงเสียดทานจลน์ $\mu_k = 0.40$ และกำหนด $g = 9.8\text{ m/s}^2$ จงหา
\begin{enumerate}
    \item สัมประสิทธิ์แรงเสียดทานสถิต $\mu_s$ ระหว่างบล็อกไม้กับพื้นเอียง
    \item ความเร่ง $a$ ของบล็อกไม้เมื่อพื้นเอียงทำมุม $45^\circ$
\end{enumerate}

\textbf{วิธีทำ:}
1. คำนวณสัมประสิทธิ์แรงเสียดทานสถิตจากมุมวิกฤต
\begin{equation}
\mu_s = \tan\alpha_c = \tan 30^\circ = \frac{1}{\sqrt{3}} \approx 0.577
\end{equation}

2. คำนวณความเร่งที่มุม $\alpha = 45^\circ$ โดยใช้สมการ \eqref{eq:kinetic_accel}
\begin{align}
a &= g(\sin 45^\circ - \mu_k \cos 45^\circ) \notag \\
&= 9.8 \left(\frac{\sqrt{2}}{2} - 0.40 \times \frac{\sqrt{2}}{2}\right) \notag \\
&= 9.8 \times \frac{\sqrt{2}}{2} (1 - 0.40) = 9.8 \times 0.7071 \times 0.60 \approx 4.16\text{ m/s}^2
\end{align}
ดังนั้น $\mu_s \approx 0.577$ และบล็อกไม้จะเคลื่อนที่ลงด้วยความเร่ง $4.16\text{ m/s}^2$
\end{workedexamplebox}

\subsection{การจำลองการถ่ายโอนพลังงานความร้อนและกฎการเย็นตัว}

นอกเหนือจากกลศาสตร์ เอนจินฟิสิกส์เชิงสะเต็มยังรองรับการจำลองอุณหพลศาสตร์ (Thermodynamics) เช่น การให้ความร้อนแก่น้ำในบีกเกอร์ด้วยขดลวดความร้อนไฟฟ้าเสมือน กำลังไฟฟ้าที่ป้อนเข้าสู่ระบบคือ $P = IV$ ซึ่งเปลี่ยนเป็นพลังงานความร้อน $Q = Pt$ ทำให้อุณหภูมิของน้ำมวล $m$ ความจุความร้อนจำเพาะ $c$ เพิ่มขึ้นตามสมการ $Q = mc\Delta T$

ในขณะเดียวกัน ความร้อนจะสูญเสียสู่สิ่งแวดล้อมตามกฎการเย็นตัวของนิวตัน (Newton's Law of Cooling) ส่งผลให้อัตราการเปลี่ยนแปลงอุณหภูมิสอดคล้องกับสมการเชิงอนุพันธ์
\begin{equation}
\frac{dT}{dt} = \frac{P - hA(T - T_{\text{env}})}{mc} \label{eq:thermo_cooling}
\end{equation}
โดยที่ $h$ คือสัมประสิทธิ์การถ่ายโอนความร้อน $A$ คือพื้นที่ผิวสัมผัส และ $T_{\text{env}}$ คืออุณหภูมิห้อง การอินทิเกรตสมการนี้ในลูปฟิสิกส์ช่วยให้ผู้เรียนเข้าใจสภาวะสมดุลความร้อน (Thermal Equilibrium) เมื่ออัตราการรับความร้อนเท่ากับอัตราการสูญเสียความร้อนพอดี ($dT/dt = 0$)

\clearpage

\section*{สรุปท้ายบทที่ 4}
\addcontentsline{toc}{section}{สรุปท้ายบทที่ 4}

เอนจินฟิสิกส์ที่มีความเที่ยงตรงสูงเป็นหัวใจของการเรียนรู้สะเต็มศึกษาในโลกเสมือนจริง การเลือกใช้ระเบียบวิธีอินทิเกรตเชิงตัวเลขที่เหมาะสม เช่น Velocity Verlet ร่วมกับการคำนวณการชนด้วยแรงดล การจำลองลูกตุ้มอย่างง่าย และการทดลองพื้นเอียงปรับมุม ช่วยสร้างสภาพแวดล้อมการทดลองที่ถูกต้องตามทฤษฎีทางวิทยาศาสตร์ เปิดโอกาสให้ผู้เรียนสามารถทดลองและสังเกตผลทั้งในเชิงกลศาสตร์และอุณหพลศาสตร์ได้อย่างแม่นยำ

\clearpage

\section*{แบบฝึกหัดท้ายบทที่ 4}
\addcontentsline{toc}{section}{แบบฝึกหัดท้ายบทที่ 4}

\begin{enumerate}
    \item จงอธิบายเหตุผลทางฟิสิกส์ว่าทำไมวิธี Explicit Euler จึงทำให้พลังงานรวมของระบบมวล-สปริงเพิ่มขึ้นอย่างต่อเนื่องเมื่อเวลาผ่านไป
    \item กำหนดให้ลูกบอลยางมวล $0.1\text{ kg}$ ตกกระทบพื้นเสมือนด้วยความเร็ว $4.0\text{ m/s}$ และกระดอนขึ้นด้วยความเร็ว $3.2\text{ m/s}$ จงคำนวณสัมประสิทธิ์การคืนตัว $e$ และขนาดของการดล $J$ ที่พื้นกระทำต่อลูกบอล
    \item เขียนโปรแกรมจำลองการสั่นแบบหน่วงวิกฤต ($\zeta = 1.0$) ด้วยวิธี Semi-Implicit Euler ในภาษา Python
    \item ลูกตุ้มอย่างง่ายมีความยาวเชือก $L = 0.98\text{ m}$ แขวนอยู่ในสนามโน้มถ่วง $g = 9.8\text{ m/s}^2$ จงคำนวณคาบการแกว่งที่มุมขนาดเล็ก และอธิบายว่าทำไมการแกว่งที่มุมกว้าง $60^\circ$ จึงต้องอาศัยการอินทิเกรตเชิงตัวเลข
    \item บล็อกโลหะวางบนพื้นเอียงเสมือนจริง เริ่มไถลลงเมื่อทำมุมเอียง $36.87^\circ$ จงหาสัมประสิทธิ์แรงเสียดทานสถิต $\mu_s$
    \item ขดลวดความร้อนไฟฟ้ากำลัง $P = 500\text{ W}$ ให้ความร้อนแก่น้ำมวล $m = 0.5\text{ kg}$ ($c = 4184\text{ J/(kg}\cdot^\circ\text{C)}$) จงคำนวณเวลาที่ใช้ในการเพิ่มอุณหภูมิจาก $25^\circ\text{C}$ เป็น $75^\circ\text{C}$ หากไม่คิดการสูญเสียความร้อนสู่สิ่งแวดล้อม
\end{enumerate}
